\section{Introduction}
\label{sec:intro}

The paper from a 10,000 foot view:

We ported the \textbf{ITrees library:}
\begin{itemize}
  \item ${}< 16$k LOC of Rocq
  \item ${}< 330$ proofs by \textbf{coinduction}
\end{itemize}
from the \code{paco} library to the \code{rocq-coinduction} library.

Why?

\textbf{Hypothesis: Using \code{rocq-coinduction} will give shorter and more
uniform proofs by coinduction than with \code{paco} (and using it requires fewer
tactics) in ITrees.}

\note{\textbf{Structure, from the talk.} The talk has two major sections:
\emph{Context}, the theoretical foundations of coinduction and the evolution of
those techniques in Rocq (\Cref{sec:background}--\Cref{sec:paco}), and
\emph{Experiment Report} (\Cref{sec:tower}--\Cref{sec:results}). This
outline follows that order. The one structural difference is that the talk's
44 context slides cannot become half the paper; see the note at the head of
\Cref{sec:background}.}

\note{\textbf{Add a contributions list at the end of this section.} CPP
reviewers look for one, and the talk does not have a slide for it. The four the
material supports: the refactoring itself; the negative result from the first
port; the rewritten library; and the three-checkpoint quantitative evaluation.}
